% =================================================================
% 1. COMPTEURS ET PACKAGES DE BASE
% =================================================================
\usepackage[dvipsnames,svgnames]{xcolor}
\usepackage[most,skins,breakable]{tcolorbox}
\usepackage[column=X]{cellspace}

% Initialisation des compteurs (partent de 0 pour l'auto-incrément)
\newcounter{DM} \newcounter{Interro} \newcounter{DS}
\newcounter{exo} \newcounter{thm} \newcounter{prp} \newcounter{dfn}

\setlength{\cellspacebottomlimit}{4pt}
\setlength{\cellspacetoplimit}{4pt}
\setlength{\parindent}{2pt}
\setlength{\tabcolsep}{0.5cm}

\pagestyle{empty} % Pas de numérotation pour les copies de devoirs

% =================================================================
% 2. BOITES TCOLORBOX (STYLE DEVOIRS)
% =================================================================

% Style générique pour les exercices de type "bordure gauche"
\tcbset{styleExo/.style={
    enhanced, breakable, frame empty, nobeforeafter,
    boxrule=0pt, titlerule=0pt,
    toprule at break=0pt, bottomrule at break=0pt,
    colback=white, sharp corners,
    borderline west={0.5mm}{0mm}{green!55!black},
    attach boxed title to top left={xshift=-3mm},
    boxed title style={size=small, colback=green!55!black, colframe=green!55!black, arc=3mm}
}}

% Exercice simple
\newtcolorbox{exercice}{styleExo, 
    title=Exercice n°\theexo, 
    code={\stepcounter{exo}}}

% Exercice avec titre personnalisé
\newtcolorbox{exerciceTitre}[1]{styleExo, 
    title=Exercice n°\theexo\ - #1, 
    code={\stepcounter{exo}}}

% Exercice noté (ex: Exercice 1 - 4 pts ou 2.5 pts)
\newtcolorbox{exerciceNote}[1]{styleExo, 
    title=Exercice n°\theexo\ \hfill (#1 pt\ifdim#1pt>1pt s\fi), 
    code={\stepcounter{exo}}}

% Théorème (cadre complet)
\newtcolorbox{theoreme}{breakable, colback=red!6!white, colframe=red!85!white,
    fonttitle=\bfseries, title=Théorème \thethm, 
    code={\stepcounter{thm}}, arc=3mm}

% Boite Titre (pour les cadres de notes ou consignes)
\newtcolorbox{boiteTitre}[2][]{
    enhanced, size=fbox, sharp corners,
    colback=white, colframe=black,
    colbacktitle=black, fonttitle=\bfseries,
    attach boxed title to top left={yshift=-3mm, xshift=2mm},
    boxed title style={size=small, left=2pt, right=2pt, sharp corners},
    title=#2, #1}

% =================================================================
% 3. EN-TÊTES (DM, DS, INTERROS)
% =================================================================

% Macro générique pour éviter la répétition (interne)
\newcommand{\layoutEntete}[3]{
    \noindent
    \begin{tabularx}{\textwidth}{@{}p{2cm} >{\centering\arraybackslash}X p{3cm}@{}}
        \textbf{\maClasse} & {\huge \textbf{#1}} & #2 \\
        & #3 & \\
    \end{tabularx}
    \vspace{0.4cm}
}

% Devoirs Maison
\newcommand{\enteteDM}[2]{\layoutEntete{Devoir maison n°#1}{Pour le #2}{}}
\newcommand{\enteteDMCorrection}[2]{\layoutEntete{Devoir maison n°#1}{Pour le #2}{\textit{Correction}}}
\newcommand{\enteteDMIndex}{\layoutEntete{Contenu des DM}{}{}}

% Interrogations
\newcommand{\enteteInterro}[3]{\layoutEntete{Interrogation n°#1}{Le #2}{Sujet #3}}
\newcommand{\enteteInterroUnique}[2]{\layoutEntete{Interrogation n°#1}{Le #2}{}}
\newcommand{\enteteInterroUniqueCorrection}[2]{\layoutEntete{Interrogation n°#1}{Le #2}{\textit{Correction}}}
\newcommand{\enteteInterroCorrection}[3]{\layoutEntete{Interrogation n°#1}{Le #2}{Sujet #3 - \textit{Correction}}}
\newcommand{\enteteInterroIndex}{\layoutEntete{Contenu des interros}{}{}}

% Devoirs Bilans (DS)
\newcommand{\enteteDS}[2]{\layoutEntete{Devoir bilan n°#1}{Le #2}{}}
\newcommand{\enteteDSCorrection}[2]{\layoutEntete{Devoir bilan n°#1}{Le #2}{\textit{Correction}}}
\newcommand{\enteteDSAB}[3]{\layoutEntete{Devoir bilan n°#1}{Le #2}{Sujet #3}}
\newcommand{\enteteDSABCorrection}[3]{\layoutEntete{Devoir bilan n°#1}{Le #2}{Sujet #3 - \textit{Correction}}}

% Divers
\newcommand{\NomPrenom}{
    \noindent \textbf{Nom Prénom :} \dotfill \medskip
}